\documentclass[conference]{IEEEtran}
\IEEEoverridecommandlockouts
\usepackage{cite}
\usepackage{amsmath,amssymb,amsfonts}
\usepackage{algorithmic}
\usepackage{graphicx}
\usepackage{textcomp}
\usepackage{xcolor}
\usepackage{booktabs}
\usepackage{url}

\def\BibTeX{{\rm B\kern-.05em{\sc i\kern-.025em b}\kern-.08em
    T\kern-.1667em\lower.7ex\hbox{E}\kern-.125emX}}

\begin{document}

\title{AI-Augmented SDN Threat Isolation: Closing the Containment Gap with Multi-Model Behavioral Telemetry and Verifiable Audit Trails}

\maketitle

\begin{abstract}
Software-Defined Networking (SDN) centralizes control-plane visibility, offering programmatic network orchestration. However, contemporary SDN threat isolation architectures remain vulnerable to high containment latencies and crude mitigation semantics, frequently resorting to unbounded source-IP blocks that inflict collateral denial-of-service on legitimate co-located tenants. In this paper, we propose an autonomous, multi-tier threat containment architecture combining edge access control, inline Suricata Deep Packet Inspection (DPI), multi-model behavioral AI, and tamper-evident cryptographic evidence sealing. Our behavioral subsystem extracts a frozen 6-feature telemetry contract from OpenFlow counters over exact measured sample intervals ($\Delta t$), explicitly detecting switch restarts and counter resets. We benchmark our pipeline on the corrected DistriNet IEEE CNS 2022 dataset suite (improved CIC-IDS-2017 and CSE-CIC-IDS-2018), resolving distribution shift through stratified sampling and sanitizing failed connection attempts. Our models achieve 99.71\% accuracy (0.9973 Macro F1) with XGBoost at 0.12~ms latency, 97.81\% accuracy (0.9792 Macro F1) with PyTorch Bi-LSTM at 0.45~ms latency, and 65.73\% anomaly F1 with an unsupervised Isolation Forest. All enforcement actions install strictly scoped 4-tuple OpenFlow rules with cryptographic cookies and expiring timeouts, verified via barrier readbacks and logged to an append-only ECDSA/SHA-256 hash-chain audit spool.
\end{abstract}

\begin{IEEEkeywords}
Software-Defined Networking, OpenFlow, Intrusion Detection, XGBoost, Bi-LSTM, Isolation Forest, Suricata, Network Security.
\end{IEEEkeywords}

\section{Introduction}
Software-Defined Networking (SDN) decouples network forwarding hardware from centralized control logic, providing global topology awareness and programmatic flow steering. Despite these architectural advantages, practical SDN threat mitigation faces two major operational challenges:

\begin{enumerate}
    \item \textbf{Containment Latency Gap:} Asynchronous counter polling introduces multi-second delays before anomalous rate spikes are detected, allowing volumetric attacks to reach line-rate saturation before remediation.
    \item \textbf{Coarse-Grained Over-Enforcement:} Traditional controllers often install broad source-IP drop rules ($nw\_src = IP$). In shared cloud subnets or environments behind Network Address Translation (NAT), this severs connectivity for thousands of benign tenants.
\end{enumerate}

To overcome these limitations, we present an end-to-end threat containment pipeline that decouples threat detection into two complementary paths (fast signature inspection and mid-flow behavioral telemetry) mediated by a risk policy engine enforcing scoped, expiring OpenFlow rules verified via barrier readbacks.

\section{System Architecture}
Our architecture coordinates four distinct planes:
\begin{itemize}
    \item \textbf{Access Gateway:} Default-deny stateful iptables boundary filtering unauthorized ports (e.g., SSH TCP 22) prior to virtual switch ingress.
    \item \textbf{Inline Suricata IPS:} AF\_PACKET bridge inspecting allowed flows (HTTPS TCP 443) and emitting structured EVE JSON alerts for exploit signatures.
    \item \textbf{Ryu OpenFlow Controller:} Operates an adaptive telemetry poller (2.0~s baseline, accelerating to 500~ms upon anomaly detection) and installs expiring mitigation rules.
    \item \textbf{Cryptographic Sealer:} Independently commits all enforcement lifecycle transitions to an append-only SHA-256 hash chain signed via ECDSA (P-256) and encrypted with AES-256-GCM.
\end{itemize}

\section{Mathematical Telemetry Formulation}
To ensure decision-time validity and prevent future-data leakage, our feature extractor derives rate metrics exclusively from active flow counter differences over measured wall-clock elapsed time $\Delta t = t_k - t_{k-1}$:

\begin{equation}
R_p = \frac{P_k - P_{k-1}}{t_k - t_{k-1}} \quad [\text{packets/sec}]
\end{equation}

\begin{equation}
R_b = \frac{B_k - B_{k-1}}{t_k - t_{k-1}} \quad [\text{bytes/sec}]
\end{equation}

\begin{equation}
S_{\text{avg}} = \frac{B_k - B_{k-1}}{P_k - P_{k-1}} \quad [\text{bytes/packet}]
\end{equation}

When switch reboots or rule evictions cause counter inversion ($P_k < P_{k-1}$ or $B_k < B_{k-1}$), the extractor asserts $\text{reset\_detected} = \text{True}$ and sets $\text{is\_valid} = \text{False}$, preventing corrupted negative rates from reaching model inference.

\section{Experimental Evaluation}

\subsection{Benchmark Dataset}
We evaluate our behavioral models using the IEEE CNS 2022 corrected NIDS dataset from DistriNet (KU Leuven). We sample a balanced benchmark of 100,000 flows across four classes: Benign, DoS/DDoS, PortScan/Infiltration, and BruteForce/Web Attacks. Failed connection attempts labeled with ``- Attempted'' are sanitized to benign per DistriNet guidelines.

\subsection{Empirical Results}
Table~\ref{tab:results} summarizes the performance across our multi-model suite evaluated on Google Colab hardware.

\begin{table}[htbp]
\caption{Multi-Model Performance on IEEE CNS 2022 Benchmark}
\label{tab:results}
\centering
\begin{tabular}{lcccc}
\toprule
\textbf{Architecture} & \textbf{Accuracy} & \textbf{Macro F1} & \textbf{Train Time} & \textbf{Latency} \\
\midrule
XGBoost & \textbf{99.71\%} & \textbf{0.9973} & 3.48~s & 0.12~ms \\
PyTorch Bi-LSTM & 97.81\% & 0.9792 & 41.57~s & 0.45~ms \\
Isolation Forest & N/A (Unsup) & 0.6573 & 0.37~s & 0.08~ms \\
\bottomrule
\end{tabular}
\end{table}

\subsection{Explainability via SHAP}
Shapley additive explanations (SHAP) reveal that $R_p$ (packet rate) and $S_{\text{avg}}$ (average packet size) serve as the primary discriminators separating volumetric flood attacks from low-rate interactive traffic and scanning probes.

\section{Conclusion}
This paper demonstrates that combining low-overhead rate telemetry, multi-model AI corroboration, and scoped, expiring OpenFlow mitigation closes the containment latency gap without inducing collateral denial-of-service. Future work will investigate hardware TCAM offloading on physical P4 switches.

\bibliographystyle{IEEEtran}
\bibliography{references}

\end{document}
